% ===================================================================
\section{Evolution: history, Darwin's logic, the synthesis (EE08)}
% ===================================================================

\subsection{What evolutionary biology must explain}
An alien visiting Earth would note four things: (1) \textbf{tremendous diversity} (Western Ghats: $\sim$7402
flowering plants, $\sim$508 birds, $\sim$6000 insects\dots); (2) \textbf{variation within groups} (skin colour, dog
breeds); (3) \textbf{complexity} (organ systems; \emph{E.\ coli} with 4402 predicted proteins); (4)
\textbf{environment-specific features} (puffins, gills). They come in two pairs: diversity and variation ask
\emph{where new forms come from}; complexity and fit ask \emph{how functional organisms arise at all}. The second pair was
the strongest argument for a designer, so a theory explaining only the first pair would not have been enough.

\subsection{The evolution of evolutionary thought}
Three positions explained diversity and complexity before Darwin: \textbf{(a)} a creator made all species,
\textbf{(b)} no species ever went extinct, \textbf{(c)} no species has changed since its origin.

\begin{center}\small
\begin{tabularx}{\textwidth}{lX}
\toprule
who & idea and effect\\\midrule
Indian scriptures & Rigveda (Hiranyagarbha, Purusha, \emph{Nasadiya Sukta}: ``no one can know''), Puranas: a creator, cycles of forms\\
Plato, Aristotle & \emph{scala naturae}, Great Chain of Being: fixed order\\
Linnaeus & taxonomy; ``as many species as the Infinite Being produced'': no extinction, no change\\
Willis, Hooke & comparative anatomy; intricacy even in a dronefly's head\\
Paley & natural theology; the \textbf{watchmaker}: complexity requires a designer\\
Steno & ``tonguestones'' are shark teeth; seashells on mountains: \textbf{the Earth changes}\\
Cuvier & mammoths: \textbf{species go extinct}, so (b) falls; but he \emph{kept} (c), no change of species\\
Taoists, Buffon & constant change; quadrupeds from \textbf{38 ancestral forms}: species can change\\\bottomrule
\end{tabularx}
\end{center}

\begin{chain}[: how the history adds up]
The first crack came from rocks, not biology: Steno showed fossils are remains of living things and that the Earth
changes. Cuvier showed some fossil species exist nowhere now, so (b) falls. If species vanish yet the world is still full
of them, new species must appear, which presses on (c); Buffon's 38 forms guessed in that direction. But none supplied a
\emph{mechanism} that also explained complexity and fit, so design (a) remained the only account of the second pair of
observations. That is the gap Lamarck and Darwin tried to fill.
\end{chain}

\subsection{What Lamarck actually said}
The caricature is ``striving produces heritable change'' (the giraffe). Lamarck (1744--1829) gave the \textbf{first
coherent theory of change in species}, in four parts:
\begin{enumerate}
\item life is continuously generated from non-living matter (spontaneous generation);
\item a \textbf{master plan} drives lineages toward greater complexity;
\item a \textbf{complexifying force} moves lineages up the plan (explains diversity and graded similarity: lineages
started at different times are at different rungs), and an \textbf{adaptive force} fits organisms to local conditions;
\item the adaptive force acts by \textbf{use and disuse} and \textbf{inheritance of acquired characters} (explains variation
within groups and adaptation).
\end{enumerate}
The theory covered all four observations. Inheritance of acquired characters was a small part and \textbf{not his idea}.
It was not accepted because of Cuvier's opposition and the lack of evidence for inheritance of acquired traits.

\begin{trap}
``Lamarck's theory = inheritance of acquired characters'' is the caricature. And the modern epigenetic cases do
\emph{not} show ``Lamarck was right'': his theory was a master plan plus a complexifying force.
\end{trap}

\subsection{Darwin's logic}
\begin{chain}[: observations $\to$ inferences, with what each borrows]
\begin{enumerate}
\item \textbf{Obs 1} overproduction (from \textbf{Malthus}: unchecked numbers grow geometrically); \textbf{Obs 2--3}
resources and population sizes stay roughly constant. These cannot all hold unless most offspring die before
reproducing, and limited resources decide who survives: \textbf{Inf 1, struggle for existence} (unequal success in a
crowded world, not literal fighting; from \textbf{Adam Smith}, the idea that competition among individuals each acting
for itself can produce order without a designer).
\item \textbf{Obs 4} individuals differ; then survival is not random with respect to those differences. \textbf{Obs 5}
offspring resemble parents, so the survivors' advantage is passed on: \textbf{Inf 2, natural selection} (also
\textbf{Obs 6}: seen in artificial selection, e.g.\ pigeons, and in the wild). Both are needed: without variation,
deaths are random with respect to traits; without heritability, each generation's advantage dies with it.
\item Selection repeated every generation, with the environment doing the ``choosing'', accumulates small advantages:
\textbf{Inf 3, adaptation}. This answers Paley: fit is what a blind, repeated filter produces.
\item Different places select in different directions; with \textbf{Lyell}'s deep time (uniformitarianism), tiny steps
add up to ``very different'' descendants: \textbf{Inf 4, new species and groups}.
\item Run Inf 4 backwards and lineages converge: all living and extinct species are branches of one tree from one or a few
ancestors: \textbf{Inf 5}, \textbf{descent with modification}.
\end{enumerate}
\end{chain}
It explains: diversity (branching); adaptations; \textbf{homology} (same bones in arm, wing, flipper: inherited and
modified); fossils resembling living species; similar species living close together; distribution breaks such as the
\textbf{Wallace line}. The lecture's caveat: this is how Darwin framed it, not how we understand it today; he had no theory
of where variation comes from or how it is inherited.

\subsection{Four objections}
Each attacks one link of the chain:
\begin{center}\small
\begin{tabularx}{\textwidth}{XXX}
\toprule
objection & link attacked & answer today\\\midrule
1. What causes variation? & source of Obs 4 & mutation and recombination (genetics)\\
2. Enough time? (\textbf{Kelvin}: a few million years) & Inf 4 needs deep time & \textbf{radiometric dating}: billions of years\\
3. Gaps in the fossil record & intermediates predicted by Inf 4--5 & fossilisation is rare; many undiscovered\\
4. Mechanism of inheritance? (\textbf{Jenkin}) & Obs 5 & \textbf{Mendel}: particulate inheritance\\\bottomrule
\end{tabularx}
\end{center}
Reactions were mixed: many accepted change but not gradualism (\textbf{saltationism}); selection as \emph{the}
mechanism stayed contested; some were disturbed by the animal origin of humans.

\begin{chain}[: Jenkin's blending argument, and Mendel's answer]
Under blending, a child is the mid-parent value: $z'=\tfrac12(z_m+z_f)$. With independent parents of variance $V$,
$\mathrm{Var}(z')=\tfrac14(V+V)=V/2$. So \textbf{heritable variance halves every generation}, $V_t=V_0/2^t$
($<0.1\,\%$ left after 10 generations). Selection would quickly run out of variation to act on, and ``Darwin had no
answer''. Under Mendelian inheritance, genes pass on unblended; allele frequencies, and hence the variance they generate,
stay constant under random mating. That is the Hardy--Weinberg law (next section).
\end{chain}

\begin{worked}[: why a young Earth hurts Darwin but not Lamarck (2025 end-sem Q6)]
Darwin's mechanism is slow, gradual and undirected, so producing all of life's diversity needs vast time; a few million
years is too short. Lamarck's is \emph{directed}: the complexifying force pushes lineages up a plan, and the adaptive force
acts within each lifetime with gains inherited, so change is purposeful and fast and does not need deep time.
\end{worked}

\subsection{Mendel and the Modern Synthesis}
\textbf{Mendel}: laws proposed 1865, published 1866, rediscovered 1900; factors do not blend. This repaired objection 4
(no dilution) and, through mutation, objection 1. \textbf{Fisher, Haldane, Wright} gave the mathematics (1900--1950);
\textbf{Stebbins} (plants), \textbf{Simpson} (fossils), \textbf{Dobzhansky} (insects), \textbf{Mayr} (systematics) the
empirical work. The \textbf{Modern Synthesis} (1930s--50s; named by Julian Huxley, 1942) united genetics,
theory, experiment, palaeontology and ecology; ``the embryologists were the only people missing''.
Chief principle (Futuyma): \textbf{adaptive evolution is caused by natural selection acting on particulate (Mendelian)
genetic variation}; mutation is the \emph{raw material} of selection, not an alternative to it; mutation, gene flow,
selection and drift cause microevolution, which over long times accounts for speciation and macroevolution.

Selected principles of the synthesis (Futuyma Box 1A): phenotype $\neq$ genotype; acquired characters are not inherited*;
variation is particulate; mutation is random with respect to need; \textbf{evolution is a change in a population, not an
individual} (allele or genotype frequencies across generations); allele frequencies change randomly (drift) or
non-randomly (selection); adaptations are traits shaped by selection; populations carry much variation and can evolve fast
without new mutations; higher taxa arise by accumulation of small differences, not ``hopeful monsters''; all life is one
tree. (* qualified by later research.) Molecular biology (sequencing; Zuckerkandl and Pauling's molecular clock, 1962)
modified the synthesis but with ``no major structural changes''.

\subsection{The Extended Evolutionary Synthesis}
\begin{center}\small
\begin{tabularx}{\textwidth}{XXX}
\toprule
MS concept $\to$ implication & case study & EES concept $\to$ implication\\\midrule
features chiefly set by genes $\to$ environment only an agent of selection & rat maternal care: high licking/grooming
\textbf{demethylates} the hippocampal glucocorticoid-receptor (GR) promoter in week 1 $\to$ more GR $\to$ calmer adults;
cross-fostering reverses it (not genetic). Humans: Romanian orphans; abused children show less GR mRNA, more methylation &
features from gene $\times$ epigenome $\to$ environment can directly alter the epigenome\\
acquired traits not inherited (Weismann) & \textbf{vinclozolin}: pregnant rat exposed, male progeny diseased for
$\ge$4 generations; \textbf{toadflax} \emph{Linaria}: the peloric form has no DNA difference, its \emph{Lcyc} gene is
methylated, stably inherited & some acquired traits can be inherited $\to$ adaptive evolution at genome + epigenome level\\
genetic information passes only within a species $\to$ lineages well defined & \emph{Elysia chlorotica} keeps algal plastids;
claimed algal genes in its genome; \textbf{recent evidence: no HGT} & information can cross species $\to$ lineages not
well defined\\\bottomrule
\end{tabularx}
\end{center}
Epigenetic and horizontal changes are faster, more dynamic and can occur in many individuals at once; they add
variation and complicate inheritance (drug-resistant bacteria, invaders). But no observation contradicts the
synthesis: the EES is an \textbf{extension, not a refutation} (Pigliucci; slide 58). Cite \emph{Elysia} as ``claimed,
then not supported''.

\subsection{Proximate and ultimate, and Dobzhansky}
Analogy 1: a death has a \emph{cause} (bullet, knife) and a \emph{motive}. Biology has proximate reasons (mechanism:
molecular biology, physiology, genetics) and ultimate reasons (why the trait exists: evolution). So ``isn't why
organisms are the way they are just molecular biology?'' gets ``yes and no''. Analogy 2: without evolution, biological
facts are isolated curiosities, ``stamps in a philately book''; evolution is the binding thread. ``Nothing in biology
makes sense except in the light of evolution'' (Dobzhansky 1973); no known observation is inconsistent with evolution.

\begin{statements}
\begin{enumerate}[label=\arabic*.]
\item ``Cuvier's work on extinction supported the idea that species change over time.''
\F He established extinction but rejected change of species.
\item ``Darwin's struggle for existence follows from overproduction alone.''
\F It needs overproduction \emph{and} roughly constant resources and population size; only together do they force most
offspring to die.
\item ``Natural selection requires variation, but heritability is optional, since survivors still differ from non-survivors.''
\F Without heritability the advantage is not passed on; nothing changes across generations.
\item ``Under blending inheritance, selection would stall because heritable variance halves each generation.''
\T Jenkin's objection; $V_t=V_0/2^t$.
\item ``Kelvin's age of the Earth was a problem for Darwin's mechanism because it is slow and undirected.''
\T Inf 4 needs deep time.
\item ``Mendel's work was central to Darwin's thinking.''
\F Published 1866, ignored until 1900; Darwin had no theory of inheritance.
\item ``The Modern Synthesis treats mutation as an alternative mechanism to selection.''
\F Mutation is the raw material of selection.
\item ``Transgenerational epigenetic inheritance in toadflax proves Lamarck's theory.''
\F It shows some acquired states are inherited; Lamarck's theory was a master plan with a complexifying force.
\item ``The EES replaces the Modern Synthesis.''
\F It extends its scope; nothing contradicts the MS.
\item ``The \emph{Elysia} sea slug is a confirmed case of animal HGT.''
\F The slide notes recent evidence that no HGT happened.
\end{enumerate}
\end{statements}

\begin{longans}[: ``Trace the development of evolutionary thought from fixity to the Modern Synthesis'' / ``Is an EES needed?'']
(1) Four observations to explain; three old positions. (2) Steno, Cuvier, Buffon removing (b) and doubting (c); Paley's
design argument. (3) Lamarck's four-part theory and why it failed. (4) Darwin's chain with Malthus, Smith, Lyell.
(5) Four objections; Jenkin's arithmetic; Mendel; Fisher--Haldane--Wright; the MS principle. (6) EES: three challenged
concepts with cases; strength of claims (extension, \emph{Elysia} retracted). (7) Proximate/ultimate and Dobzhansky.
\end{longans}
